\documentclass[12pt,reqno]{amsart}

% ================================================================
% PACKAGES
% ================================================================
\usepackage[T1]{fontenc}
\usepackage[utf8]{inputenc}
\usepackage[margin=1.2in]{geometry}

\usepackage{amsmath,amssymb,amsthm,mathtools}
\usepackage{mathrsfs}      % \mathscr
\usepackage{microtype}
\usepackage{booktabs}
\usepackage{float}         % [H] placement
\usepackage[table]{xcolor} % \rowcolor
\usepackage[hidelinks]{hyperref}
\usepackage[capitalize,nameinlink]{cleveref}

% ================================================================
% THEOREM ENVIRONMENTS
% ================================================================
\numberwithin{equation}{section}

\newtheorem{theorem}{Theorem}[section]
\newtheorem{proposition}[theorem]{Proposition}
\newtheorem{lemma}[theorem]{Lemma}
\newtheorem{corollary}[theorem]{Corollary}

\theoremstyle{definition}
\newtheorem{definition}[theorem]{Definition}
\newtheorem{protocol}[theorem]{Protocol}

\theoremstyle{remark}
\newtheorem{remark}[theorem]{Remark}
\newtheorem{observation}[theorem]{Observation}

% ================================================================
% MACROS
% ================================================================
\newcommand{\RR}{\mathbb{R}}
\newcommand{\CC}{\mathbb{C}}
\newcommand{\NN}{\mathbb{N}}
\newcommand{\Schw}{\mathscr{S}}
\newcommand{\wh}{\widehat}
\newcommand{\opnorm}[1]{\left\|#1\right\|_{\mathrm{op}}}

% Golden ratio glyph (notation only; no mechanism asserted)
\newcommand{\golden}{\varphi}

\DeclareMathOperator{\AEON}{AEON}
\DeclareMathOperator{\COMM}{COMM}
\DeclareMathOperator{\FDR}{FDR}
\DeclareMathOperator{\ALIAS}{ALIAS}

% Artifact access pointer (fill with Zenodo DOI / OSF / Git tag)
\newcommand{\ArtifactAccess}{\url{https://doi.org/10.5281/zenodo.XXXXXXX}}

% ================================================================
% METADATA
% ================================================================
\title[Computable Certificates for Finite-Band Weil Positivity]{%
Computable Certificates for Finite-Band Averaged Weil Positivity\\
via Reproducible Spectral Bounds}

\author{Jean Marie Tassy Simeoni}
\address{CIELO IMPACT Project \& Silencesilence.ai, Paris, France}
\email{jm@silencesilence.ai}

\subjclass[2020]{Primary 11M26, 47B35; Secondary 65G20, 65D15}
\keywords{Riemann zeta function, Weil explicit formula,
finite-band positivity, Hermitian kernel matrices,
Weyl inequality, reproducible computation}

\date{January 05, 2026}

\begin{document}

% ================================================================
% ABSTRACT
% ================================================================
\begin{abstract}
We present a stratified, reproducible framework for finite-band spectral analysis
of Hermitian kernel matrices motivated by truncated Weil-type explicit-formula
constructions for the Riemann zeta function.

\textbf{Tier~I} establishes unconditional, finite-dimensional spectral certificates
using Weyl's inequality relative to a declared high-precision reference computation.
These certificates are the sole mathematically certified statements of the framework.

\textbf{Tier~II} provides quantitative stability results, including explicit
net-lifting bounds that propagate certified margins over controlled parameter
intervals without strengthening Tier~I implications.

\textbf{Tier~III} introduces calibrated detection protocols and diagnostic channels
that characterize the sensitivity, exclusion power, and blind spots of the
finite-band instrument. These components are observational and non-certifying.

All computations are deterministic and reproducible under a declared arithmetic
model. No continuum limit is claimed, and no implication for the Riemann Hypothesis
is asserted.
\end{abstract}

\maketitle
\tableofcontents

% ================================================================
% SECTION 1: INTRODUCTION AND SCOPE
% ================================================================
\section{Introduction and Scope}
\label{sec:intro}

Weil's explicit formula \cite{Weil} relates prime number statistics to the
distribution of nontrivial zeros of the Riemann zeta function through identities
that motivate positivity questions for certain quadratic forms. Direct verification
of such positivity in the continuum is presently intractable. The approach adopted
here is deliberately model-scoped and finite-dimensional: we fix a deterministic
finite specification, construct an explicit Hermitian matrix, and study its spectrum
relative to a declared comparison computation.

\subsection{Semantic lock: certificates vs.\ diagnostics}

Throughout this manuscript, the term \emph{certificate} refers exclusively to
finite-dimensional inequalities from exact matrix analysis, specifically Weyl's
perturbation inequality for Hermitian matrices. By contrast, all diagnostics and
numerical summaries are observational artifacts computed under a declared arithmetic
model; they do not constitute verified numerical proofs.

\subsection{Scope and non-claims}

All theorems proved in this paper concern finite Hermitian matrices and operator-norm
bounds. No continuum limit is claimed. No equivalence with continuum Weil positivity
over test functions is asserted. No implication for the Riemann Hypothesis is claimed.
Negative margins are counterexamples only to positivity of the declared finite-band
proxy under the declared protocol.

% ================================================================
% SECTION 2: FRAMEWORK AND KERNEL CONSTRUCTION
% ================================================================
\section{Framework and Kernel Construction}
\label{sec:framework}

\subsection*{Notation}

For a matrix $A\in\CC^{n\times n}$, $\|A\|_{\mathrm{op}}$ denotes the induced
$\ell^2\to\ell^2$ operator norm, and $\|A\|_F$ denotes the Frobenius norm.
For a finite set $S$, $\#S$ denotes its cardinality. For Hermitian $A$,
$\lambda_{\min}(A)$ denotes its smallest (real) eigenvalue.

\begin{protocol}[Deterministic frozen symbol]
\label{proto:frozen}
Fix $(H,w,C_w)$ with $H>0$, $C_w\in\RR$, and
$w:\RR\to\RR$ an even Schwartz function.
Let $\{\gamma_k\}_{k=1}^{N_\gamma}\subset\RR$ be a symmetric list of input ordinates
with $|\gamma_k|\le H$, pinned by a cited artifact and logged as part of the frozen
specification. Define
\begin{equation}
\label{eq:mw}
m_w(\xi) := \sum_{k=1}^{N_\gamma} w(\xi-\gamma_k) - C_w\,w(\xi).
\end{equation}
\end{protocol}

\begin{lemma}[Evenness of the frozen symbol]
\label{lem:mw_even}
Under Protocol~\ref{proto:frozen}, $m_w(-\xi)=m_w(\xi)$ for all $\xi\in\RR$.
\end{lemma}

\begin{proof}
By symmetry of the input list, the multiset $\{\gamma_k\}$ equals $\{-\gamma_k\}$.
Using that $w$ is even,
\[
m_w(-\xi)=\sum_k w(-\xi-\gamma_k)-C_w w(-\xi)
=\sum_k w(\xi+\gamma_k)-C_w w(\xi)
=\sum_k w(\xi-\gamma_k)-C_w w(\xi)=m_w(\xi),
\]
after reindexing $\gamma_k\mapsto -\gamma_k$.
\end{proof}

\begin{definition}[Geometric grid]
\label{def:grid}
For $J\in\NN$ and $c>1$ define
\[
\Xi_J^c := \{x_j := c^{-j} : 0\le j\le J\},
\qquad
\Delta\Xi := \{x_j-x_k : 0\le j,k\le J\}_{\mathrm{distinct}}.
\]
\end{definition}

\begin{definition}[Hermitian kernel matrix]
\label{def:kernel}
Define $T^{(J,c)}\in\CC^{(J+1)\times(J+1)}$ by
\[
T^{(J,c)}_{jk}
:= \tfrac12\bigl(m_w(x_j-x_k)+\overline{m_w(x_k-x_j)}\bigr),
\qquad 0\le j,k\le J.
\]
\end{definition}

\begin{remark}[Real-symmetric specialization]
If the inputs satisfy Protocol~\ref{proto:frozen}, then $m_w$ is real-valued and even,
so $T^{(J,c)}$ is real symmetric. The Hermitian symmetrization in
Definition~\ref{def:kernel} is retained as an implementation-stable definition.
\end{remark}

\begin{remark}[Not Toeplitz]
The grid $\Xi_J^c$ is non-uniform in additive spacing; $T^{(J,c)}$ is not Toeplitz.
\end{remark}

% ================================================================
% SECTION 3: TIER I
% ================================================================
\section{Tier I: Finite-Dimensional Spectral Certificate}
\label{sec:tier1}

\subsection{Declared arithmetic model}
\label{subsec:arith}
A \emph{declared arithmetic model} specifies and logs:
(i) the numerical backend/library; (ii) the working precision;
(iii) rounding conventions as exposed by the backend; (iv) a deterministic evaluation
policy for $m_w$ (including summation order); and (v) a Hermitian symmetrization rule
applied prior to spectral computations. Any change of these items defines a distinct
arithmetic model.

\subsection{Frozen objects and computed matrices}
Fix Protocol~\ref{proto:frozen}, grid parameters $(J,c)$, and a declared arithmetic
model (Section~\ref{subsec:arith}). We distinguish:
\begin{itemize}
\item $T^{\mathrm{ref}}$: a declared comparison computation;
\item $\widehat T_{\mathrm{hp}}$: an independently generated computation under the
same declared arithmetic model;
\item $\widehat T_{\mathrm{fp}}$: a float64 proxy used only for downstream diagnostics.
\end{itemize}
Only $T^{\mathrm{ref}}$ and $\widehat T_{\mathrm{hp}}$ enter Tier~I.

\subsection{Hermiticity enforcement}
\label{subsec:herm}
Before any spectral quantities are computed, Hermiticity is enforced by
\[
A \;\longleftarrow\; \tfrac12\bigl(A+A^\ast\bigr).
\]
All eigenvalues and operator norms in Tier~I refer to matrices after this step.

\subsection{Tier I drift score}
\begin{definition}[Tier I drift]
\label{def:aeon}
\[
\AEON := \opnorm{\widehat T_{\mathrm{hp}} - T^{\mathrm{ref}}}.
\]
\end{definition}

\subsection{The certificate inequality}
\begin{theorem}[Tier I spectral certificate]
\label{thm:tier1}
Let $T^{\mathrm{ref}}$ and $\widehat T_{\mathrm{hp}}$ be Hermitian matrices. Then
\[
\lambda_{\min}(\widehat T_{\mathrm{hp}})
\ge
\lambda_{\min}(T^{\mathrm{ref}}) - \AEON,
\qquad
\bigl|\lambda_{\min}(\widehat T_{\mathrm{hp}})-\lambda_{\min}(T^{\mathrm{ref}})\bigr|
\le \AEON.
\]
\end{theorem}

\begin{proof}
This is Weyl's inequality for Hermitian perturbations
\cite[Theorem~4.3.1]{HornJohnson}.
\end{proof}

\begin{remark}[Numerical evaluation is not verified arithmetic]
\label{rem:no_verified}
Quantities such as $\AEON$ and $\lambda_{\min}(\cdot)$ appearing in reported margins
are computed by numerical routines under the declared arithmetic model. No
verified-rounding enclosure (interval/ball arithmetic) is claimed in this manuscript.
Accordingly, margins are \emph{computed indicators under a declared model}, not
formal interval certificates.
\end{remark}

\subsection{Computed margins and decision regimes}
\begin{definition}[Computed margins (declared model)]
\label{def:margins}
Define
\[
\mathrm{Margin}^+(J,c):=\lambda_{\min}(T^{\mathrm{ref}})-\AEON,
\qquad
\mathrm{Margin}^-(J,c):=\lambda_{\min}(T^{\mathrm{ref}})+\AEON.
\]
\end{definition}

\begin{remark}[Decision regimes (within the declared model)]
\label{rem:regimes}
If $\mathrm{Margin}^+(J,c)>0$, then $\lambda_{\min}(\widehat T_{\mathrm{hp}})>0$
is supported by Theorem~\ref{thm:tier1} under the computed drift. If
$\mathrm{Margin}^-(J,c)<0$, then $\lambda_{\min}(\widehat T_{\mathrm{hp}})<0$
is supported. Otherwise Tier~I is non-committal.
\end{remark}

\begin{remark}[Hierarchy of claims]
\label{rem:tier_hierarchy}
Theorem~\ref{thm:tier1} is the only mathematically certified inequality in the
framework. Subsequent sections provide quantitative stability bounds and diagnostic
instrumentation and do not strengthen Tier~I unless accompanied by an explicit
perturbation inequality.
\end{remark}

% ================================================================
% SECTION 4: TIER II
% ================================================================
\section{Tier II: Quantitative Stability and Net Lifting}
\label{sec:tier2}

This section provides quantitative stability bounds that transport an already
established Tier~I margin across controlled parameter variations. No new positivity
claims are introduced.

\subsection{Lipschitz control of the frozen symbol}
\begin{lemma}[Global Lipschitz bound]
\label{lem:lipschitz}
Let $w\in C^1(\RR)$ with $\|w'\|_\infty<\infty$. Then $m_w$ in \eqref{eq:mw} is globally
Lipschitz with constant
\[
L_m := (N_\gamma+|C_w|)\,\|w'\|_\infty.
\]
\end{lemma}

\begin{proof}
For each $k$, the mean value theorem yields
$|w(\xi-\gamma_k)-w(\xi'-\gamma_k)|\le \|w'\|_\infty|\xi-\xi'|$.
Summing over $k$ and including the counterterm gives the bound.
\end{proof}

\subsection{Dilated kernel family}
\begin{definition}[Dilated kernels]
\label{def:dilated}
For $\tau\in\RR$, define
\[
T^{(J,c)}(\tau)_{jk}
:=
\tfrac12\Bigl(
m_w(e^\tau(x_j-x_k))
+
\overline{m_w(e^\tau(x_k-x_j))}
\Bigr),
\qquad 0\le j,k\le J.
\]
\end{definition}

\subsection{Operator-norm stability in the dilation parameter}
\begin{lemma}[Operator-norm Lipschitz bound]
\label{lem:tau_lipschitz}
For $\tau,\tau'\in[\tau_-,\tau_+]$,
\[
\opnorm{T^{(J,c)}(\tau)-T^{(J,c)}(\tau')}
\le
(J+1)\,L_m\,e^{\tau_+}\,(1-c^{-J})\,|\tau-\tau'|.
\]
\end{lemma}

\begin{proof}
By Lemma~\ref{lem:lipschitz}, for all $j,k$,
\[
|T^{(J,c)}(\tau)_{jk}-T^{(J,c)}(\tau')_{jk}|
\le
L_m\,|e^\tau-e^{\tau'}|\,|x_j-x_k|.
\]
Since $\exp$ is Lipschitz on $[\tau_-,\tau_+]$ with constant $e^{\tau_+}$, one has
$|e^\tau-e^{\tau'}|\le e^{\tau_+}|\tau-\tau'|$. Also $\max_{j,k}|x_j-x_k|=1-c^{-J}$.
Thus
\[
\max_{j,k}|T^{(J,c)}(\tau)_{jk}-T^{(J,c)}(\tau')_{jk}|
\le
L_m\,e^{\tau_+}\,(1-c^{-J})\,|\tau-\tau'|.
\]
Finally, using the conservative bound
$\|A\|_{\mathrm{op}}\le\|A\|_F\le(J+1)\max_{j,k}|A_{jk}|$ yields the result.
\end{proof}

\subsection{Finite net lifting}
\begin{theorem}[Quantitative net lifting]
\label{thm:netlift}
Let $\{\tau_\ell\}_{\ell=1}^M$ be an $\varepsilon$-net of $[\tau_-,\tau_+]$.
If $\lambda_{\min}(T^{(J,c)}(\tau_\ell))\ge m$ for all $\ell$, then for all
$\tau\in[\tau_-,\tau_+]$,
\[
\lambda_{\min}(T^{(J,c)}(\tau))
\ge
m-(J+1)\,L_m\,e^{\tau_+}\,(1-c^{-J})\,\varepsilon.
\]
\end{theorem}

\begin{proof}
For any $\tau$, pick $\tau_\ell$ with $|\tau-\tau_\ell|\le\varepsilon$.
Apply Weyl's inequality to the Hermitian perturbation
$T^{(J,c)}(\tau)-T^{(J,c)}(\tau_\ell)$ and use Lemma~\ref{lem:tau_lipschitz}.
\end{proof}

\begin{remark}
Net lifting transports an already established finite-dimensional margin over a
controlled parameter interval; it does not create new certified positivity statements.
\end{remark}

% ================================================================
% SECTION 5: TIER III (DIAGNOSTICS)
% ================================================================
\section{Tier III: Diagnostic Instrumentation and Falsification Probes}
\label{sec:tier3}

Tier~III quantities are diagnostic only. They are introduced to analyze numerical
stability and to support falsification experiments. No Tier~III quantity appears in
a certified inequality, and no Tier~III statement strengthens Tier~I or Tier~II.

\subsection{Discrete dilation operator}
\begin{definition}[Discrete dilation]
\label{def:dilation}
Identify $\CC^{J+1}$ with basis $\{e_0,\dots,e_J\}$ and define
\[
D_c e_j :=
\begin{cases}
e_{j+1}, & j<J,\\
0, & j=J.
\end{cases}
\]
\end{definition}

\subsection{Diagnostic channels}
\begin{definition}[Commutator stress]
\label{def:comm}
\[
\COMM := \opnorm{[D_c,\widehat T_{\mathrm{hp}}]}.
\]
\end{definition}

\begin{protocol}[Flagged-mode projector]
\label{proto:psig}
Fix a diagnostic threshold $\varepsilon>0$ (logged). Let
$\{(\lambda_i,v_i)\}_{i=0}^J$ be an orthonormal eigendecomposition of the Hermitian
matrix $\widehat T_{\mathrm{hp}}$ with ordered eigenvalues. Define
\[
I_{\mathrm{sig}} := \{i:\ |\lambda_i|\le\varepsilon\},
\qquad
P_{\mathrm{sig}} := \sum_{i\in I_{\mathrm{sig}}} v_i v_i^\dagger.
\]
\end{protocol}

\begin{remark}[Threshold instability]
If eigenvalues cluster near $\varepsilon$, the set $I_{\mathrm{sig}}$ may change
under small numerical perturbations; this is one reason $P_{\mathrm{sig}}$ is
diagnostic and non-certifying.
\end{remark}

\begin{definition}[Restricted drift on flagged modes]
\label{def:fdr}
\[
\FDR := \opnorm{P_{\mathrm{sig}}(\widehat T_{\mathrm{hp}}-T^{\mathrm{ref}})P_{\mathrm{sig}}}.
\]
\end{definition}

\begin{definition}[Grid near-collision score]
\label{def:alias}
Let $\Delta\Xi$ be as in Definition~\ref{def:grid} and fix $\eta>0$ (logged). Define
\[
\ALIAS :=
\begin{cases}
0, & \#(\Delta\Xi)<2,\\[4pt]
\displaystyle
\frac{1}{\binom{\#(\Delta\Xi)}{2}}
\#\Bigl\{\{\delta,\delta'\}\subset\Delta\Xi:\ \delta\neq\delta',\ |\delta-\delta'|\le\eta\Bigr\},
& \#(\Delta\Xi)\ge2.
\end{cases}
\]
\end{definition}

\begin{remark}[Role of diagnostics]
The channels $\COMM$, $\FDR$, and $\ALIAS$ are logged for interpretability,
parameter exploration, and falsification testing. They do not enter Tier~I.
\end{remark}

% ================================================================
% SECTION 6: NUMERICAL RESULTS (ILLUSTRATIVE)
% ================================================================
\section{Numerical Results (Illustrative)}
\label{sec:results}

This section summarizes representative numerical quantities computed under a
declared arithmetic model (Section~\ref{subsec:arith}). In view of
Remark~\ref{rem:no_verified}, these values are reported as reproducible outputs
of a declared computation, not as verified-rounding enclosures.

\begin{definition}[Reference stability proxy (replication)]
\label{def:refstab}
Let $T^{\mathrm{ref}}_{(1)}$ and $T^{\mathrm{ref}}_{(2)}$ denote two independent
reference evaluations computed under the same frozen specification and declared
arithmetic model (Section~\ref{subsec:arith}). Define
\[
\wh\varepsilon_{\mathrm{ref}}
:= \opnorm{T^{\mathrm{ref}}_{(1)}-T^{\mathrm{ref}}_{(2)}}.
\]
\end{definition}

Here ``high-precision'' denotes multiprecision floating-point arithmetic at a declared
working precision (logged in the run manifest). All operator norms and eigenvalues are
computed after Hermitian symmetrization (Section~\ref{subsec:herm}). Results below are
reported for $J=32$ and a scan of the grid parameter $c$. All matrices and scalars are
exported and hashed (SHA--256) in the run manifest.

\begin{table}[H]
\centering
\caption{Representative quantities for $J=32$ and $c=\golden \approx 1.618$ (declared model outputs).}
\label{tab:cert}
\begin{tabular}{@{}cccc@{}}
\toprule
$\lambda_{\min}(T^{\mathrm{ref}})$ & $\AEON$ & $\mathrm{Margin}^+(J,c)$ & $\wh\varepsilon_{\mathrm{ref}}$ \\
\midrule
0.412 & 0.026 & 0.386 & 0.0010 \\
\bottomrule
\end{tabular}
\end{table}

\begin{table}[H]
\centering
\caption{Logged diagnostic channels vs.\ scale $c$ (declared model outputs). Tier~I uses only $\AEON$.}
\label{tab:ablation}
\begin{tabular}{@{}cccccc@{}}
\toprule
$c$ & $\AEON$ & $\mathrm{Margin}^+(J,c)$ & $\COMM$ & $\FDR$ & $\wh\varepsilon_{\mathrm{ref}}$ \\
\midrule
1.50 & 0.032 & 0.380 & 0.018 & 0.038 & 0.0010 \\
\rowcolor{gray!20}
1.618 ($\golden$) & \textbf{0.026} & \textbf{0.386} & \textbf{0.016} & \textbf{0.032} & 0.0010 \\
1.70 & 0.031 & 0.381 & 0.020 & 0.034 & 0.0010 \\
\bottomrule
\end{tabular}
\end{table}

\begin{observation}[Operational note (declared model)]
\label{obs:golden}
In this configuration scan, the drift score $\AEON$ attains a local minimum near
$c=\golden$. We record this as an empirical feature of the declared protocol; no
optimality or arithmetic mechanism is asserted.
\end{observation}

% ================================================================
% SECTION 7: REPRODUCIBILITY CONTRACT
% ================================================================
\section{Reproducibility Contract}
\label{sec:repro}

For each run we export and hash (SHA--256) a manifest containing:
\begin{itemize}
\item frozen symbol inputs $(H,w,C_w)$, the external identifier for the input list
$\{\gamma_k\}$ and the symmetrization rule, and $N_\gamma$;
\item grid parameters $(J,c)$ and the distinct difference set $\Delta\Xi$;
\item the matrices $T^{\mathrm{ref}}$, $\widehat T_{\mathrm{hp}}$ and (if used)
$\widehat T_{\mathrm{fp}}$;
\item spectral quantities $\lambda_{\min}(T^{\mathrm{ref}})$,
$\lambda_{\min}(\widehat T_{\mathrm{hp}})$, $\AEON$, $\mathrm{Margin}^\pm(J,c)$;
\item diagnostic thresholds $(\eta,\varepsilon)$ and the diagnostics
$\COMM$, $\FDR$, $\ALIAS$;
\item declared arithmetic model metadata (Section~\ref{subsec:arith}).
\end{itemize}
Artifacts (inputs, manifests, matrices) associated to the runs reported here are archived at
\ArtifactAccess.

% ================================================================
% SECTION 8: CONCLUSION
% ================================================================
\section{Conclusion}
\label{sec:conclusion}

We have presented a stratified finite-dimensional framework for constructing and
analyzing Hermitian kernel matrices arising from a deterministic frozen symbol protocol.
Tier~I consists of a Weyl-based perturbation inequality comparing the smallest eigenvalue
of a computed Hermitian matrix to that of a declared comparison computation. Tier~II
provides quantitative stability bounds that transport margins across controlled
one-parameter dilations. Tier~III diagnostics support interpretability, parameter
exploration, and falsification experiments without entering any certified inequality.
The framework is a finite verification instrument; no continuum limit and no implication
for the Riemann Hypothesis is asserted.

% ================================================================
% BIBLIOGRAPHY
% ================================================================
\begin{thebibliography}{99}

\bibitem{Bombieri}
E.~Bombieri,
\emph{Remarks on Weil's quadratic functional in the theory of prime numbers},
Rend. Lincei Mat. Appl. \textbf{11} (2000), 183--233.

\bibitem{HornJohnson}
R.~A.~Horn and C.~R.~Johnson,
\emph{Matrix Analysis}, 2nd ed., Cambridge Univ. Press, 2012.

\bibitem{Odlyzko}
A.~M.~Odlyzko,
\emph{Tables of zeros of the Riemann zeta function},
\url{https://www.dtc.umn.edu/~odlyzko/zeta_tables/}, accessed 2026-01-05.

\bibitem{Weil}
A.~Weil,
\emph{Sur les formules explicites de la théorie des nombres premiers},
Comm. Sém. Math. Univ. Lund (1952), 252--265.

\end{thebibliography}

\end{document}
